\documentclass[11pt,a4paper]{article}

\usepackage[utf8]{inputenc}
\usepackage[T1]{fontenc}
\usepackage[spanish]{babel}
\usepackage{geometry}
\geometry{margin=2.5cm}
\usepackage{amsmath,amssymb}
\usepackage{enumitem}
\usepackage{parskip}
\usepackage{hyperref}
\hypersetup{
    colorlinks=true,
    linkcolor=blue,
    urlcolor=blue,
    citecolor=blue
}

\title{Revisión del paper \textit{biocellai\_paper.pdf}}
\author{Revisión técnica}
\date{\today}

\begin{document}

\maketitle

\section*{Resumen breve}

El paper pregunta si el texto biológico mejora representaciones celulares a nivel de donante. Para ello compara, bajo un protocolo congelado y con separación por donante, varias arms: baselines no-texto (B0 composición, B1 pseudobulk-PCA, B2 MLP supervisado), referencias contrastivas no semánticas (T0 one-hot, T1 prototipos aleatorios), arms con texto (T2 captions de marcadores, T3 captions PubMed/RAG) y nulls (N1 etiquetas permutadas, N2 mapa caption permutado). Evalúa patología a nivel de donante en SEA-AD MTG/PFC y ROSMAP, con transferencia estricta entre cohortes. Concluye que el text grounding no mejora la generalización a nivel de donante, mientras que B2 sí transfiere estrictamente, y que una parte sustancial de la señal aparente proviene de estructura de donante.

\section*{Fortalezas}

\begin{enumerate}[label=\arabic*.]
    \item \textbf{Diseño controlado y conciencia de leakage.} El paper documenta auditorías de protocolos previos (v1) y congela contratos de datos, splits, transformaciones train-only, hashes de folds y gates pre-registrados. Esto es excelente y poco común.
    \item \textbf{Inferencia a nivel de donante.} Usa bootstrap agrupado por donante, comparaciones pareadas y permutation nulls compartidos. Esto evita pseudorreplicación y es el estándar correcto para afirmaciones donor-held-out.
    \item \textbf{Matriz de arms bien pensada.} Incluye referencias no semánticas (T0/T1) que permiten separar ``geometría contrastiva'' de ``semántica biológica''. La comparación T2/T3 vs T1 es la más limpia para evaluar si el texto aporta algo.
    \item \textbf{Resultado positivo de transferencia estricta.} B2, con predictor congelado en la cohorte fuente, supera el gate $\rho \geq 0.35$ en ambas direcciones y todos los seeds. Esto es evidencia fuerte de señal portátil.
    \item \textbf{Honestidad y reproducibilidad.} Reporta gates fallidos, límites estructurales, artefactos versionados, código y predicciones por donante. La conclusión está acotada explícitamente a este diseño.
    \item \textbf{Cuantificación del confound por estructura de donante.} Aunque discutible (ver abajo), el intento de medir cuánto de la señal aparente se debe a estructura de donante es metodológicamente valioso.
\end{enumerate}

\section*{Problemas mayores}

\subsection*{1. La comparación B2 vs T2/T3 no es de ``supervisión equivalente'' para el endpoint}

B2 tiene una \textbf{cabeza de regresión de patología a nivel donante} entrenada directamente con la variable objetivo. T2/T3 se entrenan por alineamiento contrastivo contra captions y luego se les ajusta una regresión ridge sobre embeddings congelados. Por tanto:

\begin{itemize}
    \item B2 ve la patología del donante durante el entrenamiento del encoder.
    \item T2/T3 no ven la patología del donante durante el entrenamiento del encoder; solo la ven en el readout ridge.
\end{itemize}

Esto no es supervisión equivalente para el endpoint. La comparación T2/T3 vs B2 mezcla dos cosas: (a) texto vs no-texto y (b) supervisión directa del endpoint vs supervisión indirecta/surrogate. El paper reconoce en la discusión que B2 ``optimiza directamente el endpoint'', pero el abstract y la conclusión siguen diciendo ``under matched supervision''. Eso es engañoso. La comparación justa para ``¿el texto aporta semántica?'' es T2/T3 vs T1, y ahí el texto no ayuda. La comparación T2/T3 vs B2 debería presentarse como ``texto contrastivo vs supervisión directa'', no como ``texto vs supervisión equivalente''.

\textbf{Recomendación:} añadir un brazo de texto con la misma cabeza de patología que B2 (por ejemplo, T2/T3 + regression head) o reformular todas las afirmaciones para no llamar ``matched supervision'' a B2.

\subsection*{2. El null N1 es sospechoso y necesita verificación}

El paper dice que al permutar etiquetas donante$\rightarrow$patología y reentrenar B2, el null retiene $\rho \approx 0.3$--$0.4$. Si la permutación es correcta y el modelo se reentrena desde cero con esas etiquetas permutadas, la correlación esperada con la patología verdadera en donantes held-out debería ser $\sim 0$. Que sea $0.3$--$0.4$ sugiere al menos una de estas posibilidades:

\begin{itemize}
    \item La permutación no rompe completamente la asociación (p. ej., se permuta dentro de folds, o se conserva la cabeza de tipos celulares con etiquetas verdaderas y eso filtra la patología).
    \item Hay leakage del test set o de la estructura de donante que el modelo explota de forma no esperada.
    \item El readout no se reentrena realmente con etiquetas permutadas.
    \item El null no está centrado en cero por un problema de implementación.
\end{itemize}

La explicación de que ``el encoder puede memorizar identidad de donante y la estructura de donante correlaciona con patología'' no basta: si las etiquetas están permutadas, el mapa de identidad$\rightarrow$etiqueta es aleatorio; en donantes nuevos, aplicarlo no debería predecir la patología verdadera. Este punto es central porque el paper lo usa como una de sus contribuciones principales (``quantification of a subtle confound''). Si el null es inválido, esa contribución se cae.

\textbf{Recomendación:} documentar exactamente el procedimiento de permutación, verificar que el null está centrado en cero para un predictor aleatorio, y quizás añadir un null adicional donde se permuten features de donante o se entrene con etiquetas aleatorias pero sin cabeza de tipos celulares.

\subsection*{3. Poder estadístico limitado y resultados mixtos}

Los intervalos de confianza para $\Delta\rho$ en MTG son anchos ($\sim 0.5$) y varios incluyen cero. En ROSMAP la degradación de T2/T3 vs B2 es significativa, pero en MTG no. El paper lo reconoce, pero la conclusión de ``no hay beneficio'' es más fuerte de lo que el poder permite. Con 84 donantes, efectos pequeños o moderados no pueden excluirse. La frase ``no detectable benefit'' es correcta; ``no improvement'' a secas es demasiado fuerte.

\subsection*{4. Gates pre-registrados imposibles}

El gate de marker recovery P@10 $\geq 0.6$ era estructuralmente inalcanzable porque solo $13$--$33\%$ de los genes de referencia están en el espacio de 2000 HVG, con techo máximo $0.16$--$0.38$. Esto es una falla de diseño del pre-registro: un gate debe ser factible. El paper lo reporta honestamente, pero revela que la fase de pre-registro no verificó la alcanzabilidad de todos los gates. Es una lección metodológica, no un error fatal, pero debe señalarse.

\section*{Problemas menores}

\begin{itemize}
    \item \textbf{Falta detalle sobre la generación de captions T2/T3.} No queda claro cómo se construyen exactamente los captions de marcadores, cómo se hace el RAG/PubMed, qué prompts se usan, cómo se evita circularidad, etc. Para reproducibilidad, esto es esencial.
    \item \textbf{Tabla 4 incompleta y algunas figuras con formato roto.} Es un draft, pero debe corregirse.
    \item \textbf{Endpoint ROSMAP ordinal vs CPS continuo.} Ya está declarado, pero limita comparaciones directas.
    \item \textbf{PFC y MTG comparten donantes.} Está bien marcado como robustness check, pero no debe contarse como replicación independiente.
    \item \textbf{B1 iguala a B2 en ROSMAP.} Esto sugiere que el endpoint está fuertemente impulsado por composición/pseudobulk. El paper lo discute, pero refuerza que el texto no tiene mucho margen para brillar en esta tarea.
    \item \textbf{No hay evaluación de otros endpoints.} El título y la pregunta son generales, pero la evidencia es solo patología de Alzheimer en corteza post-mortem. La conclusión está acotada, pero el título podría ser más específico.
\end{itemize}

\section*{Evaluación global}

\textbf{Fortalezas:} diseño experimental riguroso, pre-registro, inferencia a nivel de donante, transferencia estricta, reproducibilidad, honestidad con resultados negativos y fallos de gates.

\textbf{Debilidades:} la comparación central B2 vs T2/T3 no es de supervisión equivalente para el endpoint; el null N1 es sospechoso y necesita verificación; poder estadístico limitado; gates imposibles; falta detalle de captions; alcance acotado.

\textbf{Calificación orientativa:} $6.5/10$. Es un paper metodológico valioso, pero necesita revisiones mayores para sostener sus afirmaciones. Si se aclara la equivalencia de supervisión, se valida el null N1 y se reformulan las conclusiones, podría ser un buen paper en un venue de bioinformática/métodos (por ejemplo, \textit{Bioinformatics}, \textit{Genome Biology} o \textit{NeurIPS Datasets \& Benchmarks}). Tal como está, el mensaje principal es prometedor pero está sobreextendido en el abstract y la conclusión.

\textbf{Recomendación:} revisión mayor. Aclarar la diferencia entre supervisión directa e indirecta, verificar el null de etiquetas permutadas, añadir un brazo de texto con cabeza de patología, y reescribir las afirmaciones para que el alcance coincida con la evidencia.

\end{document}
